\documentclass[fleqn]{article}
\usepackage[utf8]{inputenc}
\usepackage[T2A]{fontenc}
\usepackage[russian]{babel}
\usepackage{graphicx}
\usepackage{amsmath}

\renewcommand{\arraystretch}{1.5}
\renewcommand{\baselinestretch}{2.25}
\usepackage[margin=3cm]{geometry}
\usepackage{graphicx} % основной пакет для работы с графикой

\usepackage{amssymb}
\usepackage{caption}
\setlength{\mathindent}{0pt}

\begin{document}

\title{Задачка}
\author{Цуканников Кирилл Андреевич}
\maketitle
\newpage
В 3 "А" классе учатся двоечник Витя Перестукин и 28 его одноклассников. Мария Ивановна
решила провести контрольный опрос. Оценки, полученные учениками на опросе, не зависят
друг от друга. Любой из одноклассников Вити может получить любую оценку от 2 до 5
с вероятностью 25\%. У Вити Перестукина ситуация особая: если во время урока он будет
вести себя очень плохо (вероятность чего составляет 70\%), то Мария Ивановна поставит ему
кол, не проводя опроса. Если он сможет сдержаться, то Мария Ивановна опросит и его, и в
результате он с равными вероятностями может получить 2, 3 или 4 (на пятёрку его знания не
тянут ни при каких обстоятельствах). Пусть X - количество двоек, а Y - количество пятёрок
в 3 "А". Найдите дисперсии X и Y , а также ковариацию Х и Y. \\
\\
1. Сначало буду искать дисперсию количества двоек (X). \\
а) Среди 28 одноклассников \\
Проводится 1 испытание, которое может быть удачным (в нашем случае получение оценки 2) и неудачным (другие оценки). То есть получаем распределение Бернулли\\
Мат ожидание Бернулливского распределения: $MX = 0q + 1p = p$ \\
$ MX^2 = 0^2q + 1^2p = p $ \\
Дисперсия Бернулливского распределения: $DX = MX^2 - (MX)^2 = p - p^2 =  p(1-p)$ \\
Вероятность получить оценку 2 одноклассникв = 0.25. Тогда   $DX_1 = \frac{1}{4} (1-\frac{1}{4}) = \frac{3}{16}$ \\
б) Теперь считаю дисперсию для Вити.\\
У Вити 2 пути: c 0.7 ведет себя как обычно и с вероятностью 1 получает оценку 2. Либо ведет себя лучше с 0.3 и лишь с вероятностью $ \frac{1}{3} $ получает оценку 2\\
Тогда общая вероятность получить оценку 2 будет: $ 0.7 * 1 + 0.3 * \frac{1}{3} = 0.8 $ \\
Дисперсия: $ DX_2 = 0.8(1 - 0.8) = 0.16 $ \\
в) Теперь считаю общую дисперсию для класса. \\
$ DX = 28 * DX_1 + DX_2 = 28 * \frac{3}{16} + 0.16 = 5.41 $ \\
\\
2.Теперь ищу дисперсию количества пятерок (Y) \\
а) Вероятность получить 5 для одноклассников также 0.25, такое же распределение значит и дисперсия $DX_1 = DY_1 = \frac{3}{16}$ \\
б) Витя ни при каких обстоятельств не получит оценку 5. $DY_2 = 0$ \\
в) Общая дисперсия для оценки 5: $ DY = 28 * DY_1 + DY_2 = 28 * \frac{3}{16} + 0 = 5.25$\\
\\
3. Ищу ковариацию. \\
X - общее число двоек, Y - общее число пятерок \\
$ X = \displaystyle \sum_{i=1}^{29} A_i \quad Y = \displaystyle \sum_{i=1}^{29} B_i$ \\
Поищу мат. ожидание: \\
 $ MX = \displaystyle \sum_{i=1}^{29} p_{2, i} = 28 * \frac{1}{4} + 0.8 = 7.8$ (мат ожидание получения 2) \\
 $ MY = \displaystyle \sum_{i=1}^{29} p_{5, i} = 28 * \frac{1}{4} + 0 = 7$ (мат ожидание получения 5) \\
 $ Cov(X, Y) = M(XY) - MX MY $ Мы знаем MX, MY. Теперь остается найти M(XY) \\
 $ XY =  \displaystyle \sum_{i} A_i \displaystyle \sum_{j} B_j =  \displaystyle \sum_{j} \displaystyle \sum_{i} A_i B_j = \displaystyle \sum_{i} A_i B_i + \displaystyle \sum_{i \ne j} A_i B_j $ \\
 Возьму мат ожидание от этого, получаю: $ M(XY) =  \displaystyle \sum_{i} M(A_i B_j) + \displaystyle \sum_{i \ne j} M(A_i B_j) $ \\
$ \displaystyle \sum_{i} M(A_i B_j) = 0 $ так как это взаимоисключающие события, одновременно одному i-ому ученику получить оценку 2 и 5. Это невозможно. \\
$ \displaystyle \sum_{i \ne j} M(A_i B_j) = \displaystyle \sum_{i \ne j} M(A_i) M(B_j) =  \sum_{i \ne j} p_{2, i} p_{5, j} $ (так как события что i и j получат оценки независимые) \\
$  \displaystyle \sum_{i \ne j} p_{2, i} p_{5, j} =   \displaystyle \sum_{i } p_{2, i}  \displaystyle \sum_{j}  p_{5, j} -  \displaystyle \sum_{i} p_{2, i} p_{5, i} = MX MY -  \displaystyle \sum_{i} p_{2, i} p_{5, i}$ \\
Подставляю известное: $ M(XY) = 0 + (MX MY - \displaystyle \sum_{i } p_{2, i} p_{5, j}) =  MX MY - \displaystyle \sum_{i }^{29} p_{2, i} p_{5, j}$ \\
Для обычных 28 учеников $ p_2 p_5  = \frac{1}{4} \frac{1}{4} = \frac{1}{16} $. Вклад 28: $ 28 * \frac{1}{16} = 1.75 $ \\
Для Вити: $p_2 p_5 = 0.8 * 0 = 0$ Витя не может получить оценку 5 \\
Итого: $  \displaystyle \sum_{i }^{29} p_{2, i} p_{5, i} = 1.75 + 0 = 1.75 $ \\
$ MX MY = 7.8 * 7 = 54.6 $ \\
$ M(XY) = 54.6 - 1.75 = 52.85 $ \\
$ Cov(XY) = M(XY) - MX MY = 52.85 - 54.6 = -1.75 $ ковариация отрицательная, значит чем больше оценок 2, тем меньше оценок 5. \\
\\
4. Ответ: 

\begin{equation}
\boxed{
    \begin{aligned}
        DX &= 5.41 \\
        DY &= 5.25 \\
        Cov(XY) &= -1.75
    \end{aligned}
}
\end{equation}

\end{document}